\documentclass[../DoAn.tex]{subfiles}
\begin{document}
\sloppy

Chương này trình bày các công nghệ và kiến trúc được lựa chọn để xây dựng hệ
thống quản lý ký túc xá, dựa trên yêu cầu cụ thể của bài toán: xử lý phân bổ
phòng theo chu kỳ học kỳ, hỗ trợ đồng thời web lẫn di động, đảm bảo thông báo
thời gian thực và duy trì nhất quán dữ liệu khi nhiều quản trị viên thao tác
cùng lúc.

% =====================================================================
\section{Kiến trúc hệ thống}
\label{sec:architecture}
% =====================================================================

Hệ thống được tổ chức theo kiến trúc phân lớp~\cite{fowler2002patterns}: backend
chịu toàn bộ trách nhiệm xử lý nghiệp vụ và truy cập cơ sở dữ liệu, frontend
hiển thị theo mô hình server-side rendering, và một nhóm script độc lập phục vụ
các tác vụ hàng loạt như khởi tạo dữ liệu và migration.

Kiến trúc này phù hợp với đặc điểm nghiệp vụ phân bổ phòng: mỗi đợt phân bổ
cần đi qua nhiều bước theo thứ tự xác định --- kiểm tra điều kiện, tính điểm
ưu tiên, ghi kết quả --- và không thể thực thi song song trên cùng tập sinh
viên. Tách phần tính toán điểm ưu tiên ở tầng Service khỏi phần truy vấn dữ
liệu ở tầng Model giúp kiểm soát và kiểm thử từng phần độc lập.

Luồng xử lý chính: client gửi yêu cầu qua trình duyệt hoặc ứng dụng di động,
server xác thực và cấp token hoặc session, tầng Service tính điểm ưu tiên và
ghi kết quả vào cơ sở dữ liệu, sau đó phản hồi về client. Thông báo thời gian
thực được đẩy đến sinh viên qua Socket.IO~\cite{socketio2024docs} ngay sau khi
kết quả phân bổ được ghi nhận. Các tác vụ nặng như seed dữ liệu được tách hoàn
toàn ra script riêng, không ảnh hưởng đến hiệu năng API chính.

Một điểm quan trọng trong thiết kế là hệ thống lưu lại snapshot chính sách
phân bổ tại thời điểm sinh viên nộp đơn, đảm bảo khi ban quản lý thay đổi
tiêu chí sau đó, các đơn đã nộp vẫn được xét theo chính sách cũ --- đảm bảo
tính công bằng và khả năng truy vết lý do phân bổ.

% =====================================================================
\section{Công nghệ backend}
\label{sec:backend}
% =====================================================================

\subsection*{Runtime và framework}

Backend được xây dựng trên Node.js~\cite{nodejs_docs} phiên bản 22 LTS với
Express.js~\cite{expressjs_docs} 4.x. Node.js được chọn vì toàn bộ hệ thống
dùng JavaScript cả hai phía, giúp tái sử dụng logic và giảm chi phí chuyển đổi
ngôn ngữ. Vòng lặp sự kiện đơn luồng của Node.js phù hợp với đặc trưng
I/O-bound của hệ thống --- phần lớn tải đến từ truy vấn cơ sở dữ liệu và giao
tiếp mạng. Express.js được chọn vì tính tối giản, cho phép tổ chức route và
middleware theo từng nhóm nghiệp vụ như chu kỳ phân bổ, quản lý vi phạm hay
kế thừa chính sách.

\subsection*{Cơ sở dữ liệu}

Cơ sở dữ liệu sử dụng MongoDB Atlas~\cite{mongodb2024datamodel} kết nối qua
Mongoose~\cite{mongoose_docs} 8.x. Mô hình document phù hợp với dữ liệu sinh
viên và phòng ở có cấu trúc thay đổi theo từng học kỳ --- ví dụ, tiêu chí ưu
tiên có thể bổ sung trường mới mà không cần migration phức tạp. Mongoose cung
cấp schema validation tại tầng ứng dụng, pre-save hooks và phương thức truy
vấn rõ ràng. Hệ thống định nghĩa hơn 19 schema trong thư mục
\texttt{src/schemas/}, bao gồm: chu kỳ phân bổ, nhật ký kiểm toán, vi phạm
nội quy, yêu cầu bảo trì, thông báo, khai báo ưu tiên, refresh token di động
và hàng đợi outbox.

\subsection*{Ghi log, giám sát và bảo mật}

Winston~\cite{winston_docs} ghi log dạng JSON kèm metadata đầy đủ cho mọi thao
tác quan trọng; mức log điều khiển qua biến môi trường \texttt{LOG\_LEVEL}.
Sentry~\cite{sentry_docs} (\texttt{@sentry/node} 10.x) tự động thu thập error
events và performance traces trong production. Thư viện
\texttt{@opentelemetry/api}~\cite{opentelemetry_docs} cung cấp nền tảng quan
sát theo chuẩn OpenTelemetry, cho phép gắn span và trace vào các luồng nghiệp
vụ quan trọng.

Về bảo mật tầng trung gian, Helmet~\cite{helmet_docs} 8.x thiết lập HTTP
security headers, express-rate-limit~\cite{express_rate_limit} 8.x giới hạn
số yêu cầu chống brute-force, và express-mongo-sanitize chặn NoSQL injection.
Dữ liệu đầu vào được kiểm tra bằng express-validator~\cite{express_validator_docs}
và Joi~\cite{joi_docs} tùy ngữ cảnh. Module \texttt{validateEnv.js} kiểm tra
toàn bộ biến môi trường bắt buộc khi khởi động; nếu thiếu bất kỳ secret nào,
ứng dụng thoát ngay với exit code 1.

\subsection*{Lưu trữ, xuất dữ liệu và truyền sự kiện}

Cloudinary~\cite{cloudinary_docs} kết hợp Multer~\cite{multer_docs} xử lý lưu
trữ ảnh và tài liệu, đẩy thẳng lên cloud mà không lưu tạm trên đĩa máy chủ.
ExcelJS~\cite{exceljs_docs} xuất kết quả phân bổ ra file \texttt{.xlsx};
thư viện \texttt{qrcode}~\cite{qrcode_docs} sinh mã QR thẻ cư trú phía server.

Cơ chế truyền sự kiện được thiết kế hoán đổi được giữa ba lựa chọn:
in-process cục bộ (mặc định), Redis pub/sub~\cite{redis2024pubsub} và Apache
Kafka~\cite{kafka_docs}, điều khiển qua biến môi trường
\texttt{EVENT\_TRANSPORT}. Thiết kế cắm được này cho phép nâng cấp lên Kafka
khi mở rộng quy mô mà không cần sửa mã nghiệp vụ.

% =====================================================================
\section{Công nghệ frontend}
\label{sec:frontend}
% =====================================================================

Giao diện web được xây dựng theo kiến trúc đa trang dùng EJS~\cite{ejs_docs}
làm template engine phía server, kết hợp HTML5, CSS3 thuần và
Bootstrap~\cite{bootstrap_docs} 5.x. Biểu tượng sử dụng Font~Awesome~6.x~\cite{fontawesome_docs}
nhúng qua CDN.

EJS được chọn thay vì React hay Vue vì phần lớn trang cần render dữ liệu ngay
tại server trước khi gửi về trình duyệt, đơn giản hóa đáng kể logic phía
client với các biểu mẫu đăng ký nhiều bước và trang quản trị số lượng lớn bản
ghi. Tính năng bản đồ vị trí ký túc xá dùng OpenStreetMap~\cite{openstreetmap_docs}
qua phương pháp nhúng iframe, không cần thư viện JavaScript bổ sung và tương
thích hoàn toàn với Content Security Policy của hệ thống.

Hệ thống áp dụng nguyên tắc progressive enhancement: các trang cơ bản vẫn hoạt
động khi JavaScript bị tắt, các tương tác nâng cao được bổ sung bằng các
module JavaScript nhỏ tự viết.

% =====================================================================
\section{Xác thực và bảo mật}
\label{sec:security}
% =====================================================================

Hệ thống áp dụng hai cơ chế xác thực song song. Web portal dùng session-based
authentication qua \texttt{express-session}~\cite{express_session_docs}, lưu
session ID trong cookie với thuộc tính \texttt{HttpOnly}, \texttt{SameSite=strict}
và TTL 24 giờ. Mobile API dùng JWT~\cite{rfc7519} stateless với access token
TTL 15 phút và refresh token TTL 30 ngày, ký bằng hai secret riêng, quản lý
bằng \texttt{jsonwebtoken}~\cite{jsonwebtoken_docs}.

Passport.js~\cite{passportjs_docs} 0.7.x làm tầng trừu tượng hóa xác thực.
\texttt{LocalStrategy} xử lý đăng nhập bằng mật khẩu băm bcrypt~\cite{bcrypt_npm}
6.x. \texttt{GoogleStrategy}~\cite{passport_google_docs} và
\texttt{MicrosoftStrategy}~\cite{passport_microsoft_docs} theo chuẩn
OAuth~2.0~\cite{rfc6749} cho phép đăng nhập bằng tài khoản tổ chức.

Mã QR thẻ cư trú dùng HMAC-SHA256 ký payload chứa thông tin sinh viên và
timestamp, TTL 24 giờ và rotate mỗi lần gọi API. Hệ thống phát hiện bất thường
token: khi refresh token bị dùng lại từ thiết bị khác hoặc trùng \texttt{jti},
cơ chế risk score được kích hoạt và có thể vô hiệu hóa toàn bộ phiên.

Xác thực hai yếu tố được hiện thực bằng Speakeasy~\cite{speakeasy_npm} 2.x
theo chuẩn TOTP~\cite{rfc6238}, bắt buộc với tài khoản quản trị và các thao
tác nhạy cảm.

% =====================================================================
\section{Ứng dụng di động}
\label{sec:mobile}
% =====================================================================

Ứng dụng di động được phát triển bằng React~Native~\cite{reactnative2024} 0.76
kết hợp Expo~SDK~52~\cite{expo2024router}, viết hoàn toàn bằng
TypeScript~\cite{typescript_docs}. React Native được chọn vì khả năng dùng
chung phần lớn mã nguồn cho cả iOS và Android, giảm thời gian phát triển trong
bối cảnh đồ án có giới hạn nhân lực. Expo Router~v4~\cite{expo2024router} dùng
file-based routing dựa trên cấu trúc thư mục \texttt{app/}: \texttt{(tabs)/}
cho màn hình tab chính, \texttt{(auth)/} cho luồng đăng nhập, và các thư mục
riêng cho allocation, maintenance, room detail và violations.

Về quản lý trạng thái, TanStack~Query~v5~\cite{tanstack_query_docs} đảm nhiệm
server state (fetch, cache, background refetch, optimistic updates); Zustand~v5~\cite{zustand_docs}
quản lý client state. HTTP communication dùng Axios~\cite{axios_docs} với
interceptors tự động gắn JWT và xử lý refresh khi nhận 401. Kết nối realtime
qua \texttt{socket.io-client}~\cite{socketio2024docs} cho phép nhận thông báo
phân bổ tức thì. Token được lưu an toàn qua
\texttt{expo-secure-store}~\cite{expo_secure_store_docs} (Keychain/Keystore).
Mã QR thẻ cư trú hiển thị bằng
\texttt{react-native-qrcode-svg}~\cite{rn_qrcode_docs}; Sentry React Native
thu thập crash reports và performance traces.

\subsection{Progressive Web App (PWA)}
\label{subsec:pwa}

Bên cạnh ứng dụng native, hệ thống bổ sung lớp Progressive Web
App~\cite{webdev_pwa} trên web portal, dựa trên Service Worker và Web App
Manifest~\cite{w3c_manifest}.

Service Worker (\texttt{public/sw.js}) áp dụng chiến lược cache phân biệt theo
loại tài nguyên: cache-first cho asset tĩnh, network-first cho HTML động.
Khi mất mạng, phục vụ trang dự phòng \texttt{offline.html}; các yêu cầu nhạy
cảm như đăng nhập, API và Socket.IO được cấu hình bỏ qua cache.

Web App Manifest (\texttt{public/manifest.webmanifest}) khai báo chế độ hiển
thị \texttt{standalone}, bộ icon kèm \texttt{maskable}, màu chủ đề và shortcut
truy cập nhanh. Trên Android, hệ thống bắt sự kiện
\texttt{beforeinstallprompt} để hiển thị banner gợi ý cài đặt; trên iOS Safari,
hiển thị hướng dẫn thêm vào màn hình chính thủ công.

\subsection{Đóng gói và triển khai ứng dụng}
\label{subsec:mobile_deployment}

Quá trình xây dựng và phân phối ứng dụng di động sử dụng
Expo Application Services (EAS) Build~\cite{expo_eas_build}. Khác với
Expo Go chỉ phục vụ mục đích phát triển và kiểm thử, EAS Build biên dịch
ứng dụng trên hạ tầng đám mây, sinh ra các gói cài đặt độc lập phù hợp với
từng nền tảng.

Đối với Android, hệ thống được đóng gói dưới định dạng Android Package
(\texttt{.apk}) phục vụ quá trình thử nghiệm và đánh giá nội bộ. Trong môi
trường phát hành chính thức, EAS cũng hỗ trợ sinh Android App Bundle
(\texttt{.aab}) --- định dạng được Google Play khuyến nghị nhằm tối ưu kích
thước tải xuống và phân phối theo từng thiết bị~\cite{android_app_bundle}.
Quá trình build sử dụng Gradle~\cite{gradle_docs} làm hệ thống tự động hóa
biên dịch, quản lý dependency, tối ưu hóa mã nguồn và tạo gói cài đặt cuối
cùng.

Cấu hình build được khai báo trong tệp \texttt{eas.json} theo tài liệu của
Expo~\cite{expo_eas_json}, cho phép định nghĩa nhiều profile như
\texttt{development}, \texttt{preview} và \texttt{production}. Mỗi profile
thiết lập riêng chế độ tối ưu hóa, biến môi trường và khóa ký số (signing
key), nhờ đó cùng một mã nguồn có thể được sử dụng cho nhiều giai đoạn từ
phát triển, kiểm thử đến triển khai.

Việc sử dụng EAS Build giúp loại bỏ yêu cầu cài đặt Android SDK, Gradle và
các công cụ biên dịch trên máy phát triển. Thay vào đó, toàn bộ quá trình
biên dịch được thực hiện trên hạ tầng của Expo~\cite{expo_eas_build}, đảm bảo
môi trường build đồng nhất, giảm sai khác giữa các máy phát triển và đơn giản
hóa quy trình Continuous Integration/Continuous Deployment
(CI/CD)~\cite{expo_ci_cd}.
% =====================================================================
\section{Hạ tầng và vận hành}
\label{sec:infrastructure}
% =====================================================================

\subsection*{Container hóa và cache}

Backend được đóng gói trong Docker~\cite{docker_docs} container dùng multi-stage
build với base image \texttt{node:20-alpine}: stage đầu cài dependencies với
\texttt{npm ci -{}-omit=dev}, stage hai chạy dưới non-root user \texttt{appuser}.
Docker Compose phục vụ triển khai local với hai service: \texttt{app} (backend
Node.js) và \texttt{redis} (Redis~7-alpine), khởi động tuần tự theo health
check. MongoDB Atlas kết nối qua biến môi trường \texttt{MONGODB\_URI}.

Redis~\cite{redis2024pubsub} 5.x phục vụ hai mục đích: lưu session qua
\texttt{express-session} và cung cấp Socket.IO cluster adapter
(\texttt{@socket.io/redis-adapter}~\cite{socketio2024docs}) cho broadcast nhất
quán qua nhiều instance. Khi Redis không khả dụng, hệ thống fallback về chế
độ local Socket.IO. Endpoint \texttt{GET /health} phục vụ Docker và load
balancer kiểm tra trạng thái.

\subsection*{Email, CI/CD và vận hành}

Nodemailer~\cite{nodemailer_docs} 7.x gửi email qua SMTP (cấu hình
\texttt{SMTP\_*}) phục vụ xác nhận đăng ký, thông báo phân bổ, OTP và đặt
lại mật khẩu. Kênh thông báo chính vẫn là Socket.IO realtime và thông báo
in-app.

GitHub Actions~\cite{github_actions} CI pipeline chạy hai job song song:
\texttt{backend} kiểm tra syntax JavaScript bằng \texttt{node -{}-check} và
\texttt{mobile} chạy TypeScript type check bằng \texttt{npx tsc -{}-noEmit},
kích hoạt trên mọi push và pull request. Trong production, PM2~\cite{pm2_docs}
làm process manager với chế độ cluster tận dụng nhiều nhân CPU.

% =====================================================================
\section{Công nghệ trực quan hóa 3D}
\label{sec:visualization}
% =====================================================================

Hệ thống cung cấp trực quan hóa không gian theo hai cấp độ: bản đồ địa lý ba
chiều ở cấp toàn khu và tham quan thực tế ảo ở cấp từng phòng, hiện thực bằng
CesiumJS và Three.js chạy trực tiếp trên trình duyệt qua WebGL.

\subsection{CesiumJS --- Trực quan hóa địa lý 3D}
\label{subsec:cesium}

CesiumJS~\cite{cesiumjs} là thư viện JavaScript mã nguồn mở chuyên dựng địa
cầu và bản đồ ba chiều trên nền WebGL. Hệ thống tích hợp CesiumJS vào tính
năng ``Khám phá KTX 3D'' trên trang chủ, cho phép người dùng quan sát vị trí
thực tế các khu ký túc xá trong không gian địa lý. Mỗi tòa nhà được biểu diễn
bằng marker tương tác; khi chọn, hệ thống hiển thị thẻ thông tin kèm hình ảnh,
mô tả và liên kết đến trang chi tiết. Người dùng có thể phóng to, xoay và
nghiêng góc camera mượt mà trên cả máy tính lẫn di động.

\subsection{Three.js --- Tour thực tế ảo phòng ký túc xá}
\label{subsec:threejs}

Three.js~\cite{threejs} là thư viện đồ họa 3D cấp cao trên nền WebGL, dùng để
xây dựng tính năng tham quan phòng trước khi đăng ký. Mô hình phòng ở định
dạng glTF/GLB~\cite{khronos_gltf} được nạp trực tiếp trong trình duyệt qua
\texttt{GLTFLoader}. Tính năng cung cấp hai chế độ: tham quan có hướng dẫn
di chuyển camera theo lộ trình định sẵn, và khám phá tự do cho phép điều khiển
bằng chuột hoặc cảm ứng kèm xử lý va chạm cơ bản. Mô hình được nén bằng
quantization và texture WEBP để phù hợp điều kiện kết nối di động.

Tính năng này giải quyết hạn chế thực tế của các hệ thống hiện tại: sinh viên
từ tỉnh xa không cần đến tận nơi xem phòng trước khi đăng ký, tour 3D giúp
hình dung diện tích, bố trí nội thất và điều kiện sinh hoạt ngay từ trang web.

% =====================================================================
Chương 3 đã trình bày các lựa chọn công nghệ cùng lý do cụ thể: Node.js và
MongoDB Atlas cho tính linh hoạt và hiệu năng I/O, Passport.js và JWT cho xác
thực đa lớp, React Native kết hợp PWA cho khả năng tiếp cận đa thiết bị,
Docker và GitHub Actions cho vận hành nhất quán, và CesiumJS cùng Three.js cho
trải nghiệm trực quan hóa không gian. Các quyết định này tạo nền tảng cho
chương 4 trình bày thiết kế chi tiết và kết quả kiểm thử.

\end{document}